\section{Security analysis}\label{sec:security}

This section is heuristic. We state the problems the scheme relies on,
reduce key recovery to finding $f$, and estimate the cost of every attack we
know of. Table~\ref{tab:attacks} summarises the estimates. A cost written
$2^x\cdot C$ means $2^x$ evaluations of a polynomial-time subroutine $C$,
typically a $(2,2)$-chain in dimension~2 or an isogeny chain in dimension~8.

\subsection{Problems}

\begin{assumption}[hidden-degree eigenbasis problem]\label{ass:key}
Given the public parameters and $\mathit{pk}=(\EA,P_A,Q_A)$, it is hard to
find $f$, or any secret key that decapsulates honest ciphertexts for
$\mathit{pk}$.
\end{assumption}

\begin{assumption}[masked eigen-torsion problem]\label{ass:msg}
Given $\mathit{pk}$ and $\mathit{ct}=\textsf{Enc}(\mathit{pk},s)$ for uniform
$s\in\Z/d$, it is hard to find $s$. This is the one-wayness of the
deterministic PKE inside the KEM.
\end{assumption}

\begin{lemma}[the degree is the key]\label{lem:fiskey}
Given $f$, a working secret key can be computed in polynomial time, up to a
choice among the few representations of $m=2^E-7d^2f^2$ as a sum of two
squares.
\end{lemma}

\begin{proof}
$s_0=fs_D$ is public once $f$ is known. The pair $(c,c')$ is one of the
representations $c^2+c'^2=m$; key generation finds one with the same
procedure an attacker would use, and the schedule condition and the signs
leave at most a handful of candidates. Each candidate can be tested by
decapsulating a self-made ciphertext. Neither $\alpha$, $\varepsilon$ nor
$\lambda_1,\lambda_2$ is needed.
\end{proof}

So Assumption~\ref{ass:key} is equivalent to the hardness of finding $f$. A
consequence is that $\EA$ and the eigen-structure must hide $f$; the next
subsection shows what goes wrong when they do not.

\subsection{The public-degree variant is insecure}\label{sec:walk}

The toy prototype of~\cite{EigenSetaNote} had a ``walk'' variant: $\EA$ was
the codomain of a secret isogeny $\varepsilon$ of \emph{public} degree $5^k$,
and $\tau=\varepsilon\omega\hat\varepsilon$. Its decryption key
$(c,c',s_0)$ is computed from $\Delta'=4\deg\tau-(\Tr\tau)^2=7\cdot25^k$ and
from nothing else, as the merged code shows
(\art{crypto:research/eigen_seta_prototype_20261011/eigen_seta_toy.py},
lines 133--155 and 201--202). Since $k$ is public, anyone can compute the key
and decrypt. An exploratory run, not archived, decrypted 5 of 5 toy
ciphertexts this way. Version~2 exists to close this hole: the degree of
$\tau$ is $2f^2$ with $f$ secret.

\subsection{Key recovery}\label{sec:key}

\paragraph{Exhaustive search over $f$.} An attacker can enumerate candidates
$f$, apply key generation's public filters, derive a candidate key by
Lemma~\ref{lem:fiskey}, and test it with one decapsulation. The filters
accept a fraction $\rho$ of the odd candidates in $[f_{\mathrm{lo}},2f_{\mathrm{lo}})$.
We measured $\rho$ empirically
(\art{crypto:research/eigen_seta_rust_20261011/reports/run-20261011T1955Z-verify/records/f_filter_rate.txt}):
$\log_2\rho=\RhoLOne$, $\RhoLThree$ and $\RhoLFive$ at L1, L3 and L5 (measured from 200, 100 and 100 accepted candidates, so about $\pm0.1$ to $\pm0.14$ binomial standard error). This
leaves about $2^{\FspaceLOne}$, $2^{\FspaceLThree}$ and $2^{\FspaceLFive}$
admissible degrees, an upper bound since the later rejections of key
generation are not counted. The search costs that many chain evaluations
classically and their square root with Grover's algorithm. We know of no
test that rejects a wrong $f$ more cheaply than a full chain. In the toy
experiment of~\cite[\S6]{EigenSetaNote}, 85 of 86 wrong but
pairing-consistent guesses, including guesses agreeing with the truth on
3--13 low bits, ran the whole chain and failed only at the final split
(\art{crypto:research/eigen_seta_prototype_20261011/results/}).

\paragraph{Recovering $\varepsilon$ from a path.} Let $I$ be any left
$\End(\ED)$-ideal connecting $\ED$ to $\EA$. The isogeny $\varepsilon$
corresponds to an equivalent ideal of norm $f$. There are about
$\tfrac{15}{2\pi^2}X^2$ cyclic isogenies of degree at most $X$ from $\ED$,
spread over about $p/12$ curves, so on average about $9X^2/p$ of them land
on a given curve. For $X=2f_{\mathrm{lo}}$ this is about $2^{-782}$ at
level~1. So $\varepsilon$
is, with overwhelming probability, the unique shortest isogeny from $\ED$ to
$\EA$, and a reduced basis of $I$ yields it, and $f$, in polynomial time.
Key recovery therefore costs no more than finding \emph{some} isogeny from
$\ED$ to $\EA$. That is the supersingular path or endomorphism-ring
problem~\cite{EHLMP18,Wes22}: $\tilde O(p^{1/2})$ classically~\cite{DG16}
and $\tilde O(p^{1/4})$ quantumly~\cite{BJS14}. We also list $p^{1/3}$ as a
conservative classical figure in view of recent heuristic algorithms with
large memory~\cite[\S1.1]{EigenSetaNote}. Because $p\approx d^5$, all three
are far beyond the target.

\paragraph{Meet in the middle on $\varepsilon$.} If $f$ is smooth, $\varepsilon$
factors and a meet-in-the-middle search becomes possible. Even with $f$
known and perfectly balanced, it costs at least $\sqrt f\ge\sqrt{f_{\mathrm{lo}}}$:
$2^{138.5}$, $2^{204.1}$ and $2^{267.5}$ at the three levels, which is above
$2^\lambda$, and it applies only to the fraction of keys with smooth $f$.

\paragraph{Torsion-point attacks on the public key.} The public key reveals
$\varepsilon$ on $\ED[2^E]$, with $2^{2E}\gg f$, up to the diagonal scalars
$\lambda_1,\lambda_2$. The SIDH attacks~\cite{CD23,MMPPW23,Rob23} need the
degree to build the auxiliary isogeny of degree $2^E-f$ or $N^2-f$. With $f$
unknown they reduce to the exhaustive search above. Hiding the degree as a
countermeasure to these attacks also appears in terSIDH and
binSIDH~\cite{BF23}, where the secret degree ranges over products of small
prime powers; here it ranges over an interval. Knowing $f$ would make
$\lambda_1,\lambda_2$ irrelevant: $\tau$ acts on $(P_A,Q_A)$ as
$f\cdot\mathrm{diag}(\mu,\nu)$ whatever the scaling, and the attacker could
simply decrypt.

\paragraph{Leakage through the public key.} The Weil pairing
$e(P_A,Q_A)=e(P_D,Q_D)^{\lambda_1\lambda_2f}$ is uniform among primitive
$2^E$-th roots of unity, because $\lambda_1\lambda_2$ is a uniform unit.
The eigenvalues of $\tau$ are $f\mu$ and $f\nu$, with $\mu,\nu$ public, so an
attacker knows the action of $\tau$ on $\EA[2^E]$ up to the unknown scalar
$f\bmod 2^E$. We do not know how to use this without $f$. The decryption
identity $\deg\gamma+c'^2=2^E$ involves $f^2$ through $\deg\tau$, so there is
no smaller Kani instance that tests $f$ modulo $2^j$. This is the main
structural assumption of the scheme, and it is where we would look first for
an attack.

\subsection{Message recovery}\label{sec:msg}

\paragraph{Mask guessing plus a torsion attack (binding).} The ciphertext
contains $S=a\varphi P_A$ and $T=b\varphi Q_A$, with $\varphi$ of public
degree $d$. The pairing $e(S,T)=e(P_A,Q_A)^{abd}$ reveals $ab$, but the
ratio $r=b/a$ is uniform. If $r\bmod 2^n$ were known, the attacker would hold
$\varphi$ on $\EA[2^n]$ up to the common scalar $a$. Robert's attack in
dimension~8~\cite{Rob23} then recovers $\varphi$ in polynomial time as soon
as $2^{2n}>d$. With $r$ odd, guessing it costs $2^{n-1}$ trials,
where $n=\lfloor(\log_2d)/2\rfloor+1$. That is about $2^{128}$, $2^{192}$ and
$2^{256}$ dimension-8 chains classically, and the square root of that with Grover.
FESTA's analysis~\cite{BMP23} rests on the same estimate. An attacker who
knew $\tau$ could do better: the masks act diagonally on eigenvectors of
$\theta=\varphi\tau\hat\varphi$, which is exactly what decapsulation
exploits. But knowing $\tau$ means knowing $f$.

\paragraph{Claw finding on $\varphi$.} Both ends of $\varphi$ are public and
its degree is $d$. A generic meet in the middle costs about $\sqrt d$ time
and memory classically. Tani's algorithm costs $d^{1/3}$ quantum time with
as much quantum memory, but under realistic cost models it was not better
than classical van Oorschot--Wiener search~\cite{vOW99} for SIKE~\cite{JS19}. Our $d$ is
$2^{256.8}$, $2^{385.1}$ and $2^{513.5}$, which exceeds the degrees
$2^{216}$, $2^{305}$ and $2^{372}$ of SIKEp434, p610 and p751~\cite{SIKE}.
Those were assessed against the same attack for categories 1, 3 and 5. So
claw finding is not the binding attack.

\paragraph{Eigenbases of known endomorphisms.} The Castryck--Vercauteren
attacks on FESTA and M-SIDH~\cite{CV23} need an endomorphism of $\EA$ of
small degree for which $(P_A,Q_A)$ is an eigenbasis. The only such
endomorphisms we know are those in $\Z[\tau]$, of degree at least $2f^2$,
and they are not available without $f$.

\subsection{Quantum security and NIST categories}\label{sec:nist}
NIST defines categories 1, 3 and 5 by the cost of key search on AES-128, -192
and -256~\cite{NIST16}. Quantumly, that is Grover search over $2^{128}$,
$2^{192}$ or $2^{256}$ keys with an AES circuit as the oracle. The binding
attack above has the same shape: a Grover search over about $2^{\lambda}$
candidates, but with an oracle (an isogeny chain in dimension~8 over a
1341- to 2630-bit field) that is far more expensive than an AES evaluation.
Classically it is about $2^\lambda$ such chains against $2^\lambda$ AES
evaluations. We therefore \emph{conjecture} that L1, L3 and L5 meet
categories 1, 3 and 5. Two caveats apply. The margin rests on the oracle
cost alone: the search space matches the AES key space and is not larger.
And any attack that guesses less than $\sqrt d$ of mask information would
break the claim.

\subsection{Chosen-ciphertext security and decryption failures}\label{sec:cca}
The KEM is $U^{\not\perp}[\textsf{Enc},G]$~\cite{HHK17}. In the random-oracle
model it is IND-CCA secure if $\textsf{Enc}$ is one-way
(Assumption~\ref{ass:msg}) and has negligible correctness error $\delta$;
quantum random-oracle bounds for implicit rejection also
exist~\cite{BHHHP19}. The adaptive key-recovery attacks that broke static
SIDH keys~\cite{GPST16} send malformed ciphertexts and watch the result.
Here every ciphertext that differs from $\textsf{Enc}(\mathit{pk},s)$ for the
recovered $s$ yields the pseudorandom $H_R(z,\mathit{ct})$, in constant time
(\S\ref{sec:impl}). Decapsulation must also never fail \emph{observably} on
malformed inputs. The robustness test (\S\ref{sec:results}) feeds 48
adversarial ciphertexts, including singular and special curves,
small-order points and random curves: none causes a panic, and none yields
the honest key. Degenerate formula inputs do not give an attacker a failure oracle either.
A crafted ciphertext is answered with $H_R(z,\mathit{ct})$ whether or not
an intermediate formula degenerates, because either way the re-encryption
differs. Only the failure of an \emph{honest} ciphertext would be visible,
as a mismatch with the encapsulated key. That has probability
heuristically $O(E/p)$ (\S\ref{sec:correctness}), and it did not occur in
\FailTotalTrials\ decapsulations. So decryption-failure attacks have nothing
to work with.

\subsection{Weak keys}\label{sec:weak}
\begin{itemize}
\item $\alpha\in\Z[\omega]$, i.e.\ $y=0$. The Kani factor would then be an
  $\mathcal{O}_K$-ideal isogeny, and since $h(-7)=1$ its codomain is $\ED$
  itself (or its conjugate). Key generation forces $(y_0,y_1)\ne(0,0)$.
\item $f\equiv g\pmod d$. The pairing test could not tell the two Kani
  factors apart, so key generation rejects these $f$ via $2f\not\equiv2^E$.
\item Scalars $\equiv-1$ modulo the point order. These break $y$-recovery in
  the ladder; key generation rejects them.
\item The schedule condition and the gluing test. These fix the 2-adic shape
  of $(c',\lambda_S,\lambda_T)$ for all keys. They are public by design, and
  their cost in key entropy is included in $\rho$.
\end{itemize}

\begin{table}[t]
\centering
\caption[Estimated attack costs, as $\log_2$ of the number of basic
operations]{Estimated attack costs, as $\log_2$ of the number of basic
operations. ``chain'' means one isogeny-chain evaluation. All figures are
derived from the parameters, except $\rho$, which is measured. Source of
$\rho$:
\art{crypto:research/eigen_seta_rust_20261011/reports/run-20261011T1955Z-verify/records/f_filter_rate.txt}.}
\label{tab:attacks}
\small
\begin{tabular}{llrrr}
\toprule
target & attack & L1 & L3 & L5\\
\midrule
key & search over admissible $f$, classical (chains) & \FspaceLOne & \FspaceLThree & \FspaceLFive\\
key & same with Grover (chains) & \GroverFLOne & \GroverFLThree & \GroverFLFive\\
key & path to $\EA$ then shortest ideal, classical $p^{1/2}$ / $p^{1/3}$ & 670 / 447 & 995 / 663 & 1315 / 876\\
key & path to $\EA$ then shortest ideal, quantum $p^{1/4}$ & 335 & 498 & 657\\
key & meet in the middle on $\varepsilon$ ($\sqrt f$) & 138.5 & 204.1 & 267.5\\
message & mask guess $+$ dim-8 attack, classical (chains) & 128 & 192 & 256\\
message & same with Grover (chains) & 64 & 96 & 128\\
message & claw on $\varphi$, classical $\sqrt d$ & 128.4 & 192.6 & 256.8\\
message & claw on $\varphi$, quantum $d^{1/3}$ (with $d^{1/3}$ qubits) & 85.6 & 128.4 & 171.2\\
\midrule
\multicolumn{2}{l}{target $\lambda$ (classical) / NIST category} & 128 / 1 & 192 / 3 & 256 / 5\\
\bottomrule
\end{tabular}
\end{table}
